% !TeX root = ../papers/nc-p3.tex
\section{十八次での滑らかな二次曲面への双有理核の排除}
\label{sec:cofactor-eighteen-birational-quadric}
\begin{theorem}\label{thm:cofactor-eighteen-birational-quadric}
$d=18$ で核像が滑らかな二次曲面なら、核写像は次数二である。
\end{theorem}
\begin{proof}[種数三の核切断]
双有理と仮定する。逆系の一般切断 $J$ の二次数は和六、算術種数は三または四なので、
母線を交換すれば $(2,4)$ または $(3,3)$ である。前者では $p_a(J)=3$ なので
定理\ref{thm:cofactor-sectional-graph-genus}により基点スキーム $Z$ は長さ八の局所完全交叉である。
$S_0=\mathbb P^1\times\mathbb P^1$ の第一次数二の母線を繊維とする射影を $p$、
逆束を $A$ とする。$\Gamma=\operatorname{Bl}_Z S_0$ の例外直線は入力束 $H$ で次数一を持つ。
繊維 $F$ とのスキーム交点 $Z\cap F$ は長さ二以下である。入力切断は $Z$ のイデアルを生成し、
どの繊維にも少なくとも一つの非零な二次式を制限するからである。
よって $p_* A=\mathcal O(4)^3\to p_*(\mathcal O_Z A)$ は繊維ごとに全射である。
その核 $W$ は階数三、次数四である。一般繊維の入力像は滑らかな円錐曲線
（曲面への双有理性から二対一直線にはならない）ゆえ四入力切断は $W$ を一般点で生成し、全直和次数は非負である。

$\Gamma\to\mathbb P(W)$ は相対切断の評価で定まり、相対超平面 $\xi$ は $H$ に引き戻される。
実際ブローアップの入力束の $S_0$ への直像は $I_Z A$（二生成正則列の一方程式表示からも従う）で、
$W$ の評価には既に基点のない四切断が入る。この射は一般繊維で二次埋込み、像は整な因子である。
$\xi^3=4$, $H^2=8$ より因子の類はちょうど $2\xi$。
垂直な正次元入力繊維は高々四入力点に限られる。実際もとの母線の固有変換はその滑らかな母線に同型で、
$H$ の次数は $2-\ell(Z\cap F)$、長さ八の中で次数零は高々四回である。
水平な縮約の入力点は高々二つである。$W$ の次数零の直和因子の個数を $z\le2$ とすると
$|\xi|$ の水平な縮約は自明な部分スクロール $\mathbb P^1\times\mathbb P^{z-1}$ 上の定値切断だけである
（正次数の直和因子の射影座標が一つでも非零なら底座標を復元できる）。
$2\xi$ の方程式をそこに制限すると頂点空間 $\mathbb P^{z-1}$ の一つの定数係数二次式となる。
$z=2$ で恒等零なら整な像曲面がその部分スクロールに一致してしまうので不可能である。
$z=0,1$ でも個数の上界は成立する。四切断と全切断が潰す曲線は共に $H$ の次数零の曲線だから、
四切断への射影でこの個数は増えない。

長さ八の基点を平滑化でき、$W$ の非負性から $H^1(I_Z A)=0$ で入力切断も延長する
（繊維評価の全射から $R^1p_*(I_Z A)=0$）。滑らかな八点ブローアップ上で
$L^2=2$, $LK=-4$, $HL=6$, $H^2=8$, $HK=-4$ なので
$\chi(\Gamma,L+tH)=4(t+1)^2$。従って入力有限な開集合上で $M=h_*L$ が成立し、
少なくとも $20-6=14$ 個の孤立三点繊維がある。

一般八点では $W$ に次数零の因子がない。例えば相異なる八母線上で、他方の座標の三値
$0,1,\infty$ に中心を三個、三個、二個ずつ置く。二次式の三評価は独立なので $W$ の直和次数は $1,1,2$。
一般には全相対系で滑らかな八点ブローアップが滑らかなスクロールの $2\xi$ 因子として埋まる。
各特別繊維は次数 $(1,1)$ の二直線なので相対埋込みについては
定理\ref{thm:cofactor-exclude-nineteen}の証明と同じである。因子はスクロールの二次超平面切断であり、
やはり直線を張る長さ三は一般の射影で潰れない。
従って同じ変形と同伴束積分を適用でき、一般繊維の長さは十二であって十四に足りない。
\end{proof}

\begin{proof}[種数四の核切断]
$(3,3)$ の場合、$G\to J$ の正規化方向への拡大の余長は全体で一である。
非自明な点 $q$ は基点の一つであり、一般切断を選んでこの点を固定できる
（有限候補で、各基点での等号は基点切断比による有限長の条件、または一般平面の関数体上で判定できる）。
$q$ は $J$ の特異点なので逆系の重複度 $m$ は二以上。全基点の重複度平方和は十なので $m\le3$。

$m=3$ を除く。残る基点はちょうど一単純点である。$q$ を吹いた例外直線上で先頭三次式の系の固定因子の次数は
高々一である（最終系が自由で、例外直線上で除ける基点は高々一つで単純）。
従って一般 $J$ は $q$ で相異なる三接線の三枝を持つ。
$G$ はこの三枝を別の点に分けられない。完備局所環で枝を非自明な二群に分ければ、
両群の局所環の直積に入れるだけで既に余長は二以上となる。
実際二群の曲線方程式を $u,v$ とするとその余長は $(u,v)$ の交点数で、三接線の場合は二である。
中間環で点が分かれれば冪等元を持ち、この直積を含まなければならない。
従って三枝は入力平面八次の同一点に行く。その入力点は一般平面で $C$ の外にある。
実際 $q$ を消す固定列の曲線部分は次数二以下であり、$C$ との共通点は有限である。
よって十八の二重点以外に正規化余長三以上の点があり、平面八次の正規化種数は零以下。
一方 $(3,3)$ 切断の種数は四から三重点の三だけ下がり一である
（基点解消後のBertiniと随伴による）。これは矛盾である。

$m=2$ とする。$q$ の一点ブローアップを $S_1$、$J$ の固有変換を $J'$ とする。
$J'\to J$ は余長一の部分正規化である。これは二重点を一回吹いたときの算術種数差である。
Gorenstein曲線局所環 $R$ で余長一の中間環は $\operatorname{Hom}(\mathfrak m,R)$ に等しく一意である
（定理\ref{thm:cofactor-curve-kernel-only-bezout}の証明）。
従って $J'$ と $G$ は同じ部分正規化である。
一点を吹いた積二次曲面を平面二点のブローアップとして書き、平面超平面を $h$、例外類を $e_1,e_2$ とする。
二母線類は $h-e_i$、$q$ 上の例外類は $h-e_1-e_2$ ゆえ
\[
        A'=4h-e_1-e_2,\qquad L=2h-e_1-e_2 .
\]
変換後の全ての基点で $J'=G$ 上に入力比が正則なので、前定理の二生成の議論で
残余基点 $Z'$ は局所完全交叉で長さ $A'^2-8=6$ となる。
$J'$ の上でそのイデアルは可逆で、$I_{Z'}A'$ を一般切断で割った層は
$J'$ 上の次数八の線束である。算術種数三なのでその第一コホモロジーは零、
従って $H^1(I_{Z'}A')=0$。$\Gamma'=\operatorname{Bl}_{Z'}S_1$ で $H$ は自由、
$Z'$ と入力四切断を一般の六点と四切断に変形できる。一般には平面八点の
$4h-\sum e_i$、$L=2h-e_1-e_2$ ゆえ再び $\chi(L+tH)=4(t+1)^2$ となる。

入力に潰れる正次元繊維は高々四つである。$Z'$ 上の例外直線は入力次数一であり、
$S_1$ 上の二例外曲線の固有変換は高々二本である。それ以外の縮約曲線は平面上に正次数の像を持つ。
$H-h$ の有効系の全切断はこの各曲線上で負次数なので零になる。
この系には少なくとも二独立な平面三次式がある。二点を通る三次式八次元に
$S_1$ 上の長さ六の $Z'$ で消える条件を課すだけだからである。
従って固定される平面曲線の次数の和は二以下となる。以上で個数の上界を得る。
有限な開集合での押出し同型から孤立三点繊維は十六個以上。
これらが変形する一般平面八点モデルの長さは定理\ref{thm:cofactor-eighteen-smooth-cubic}で十三であった。
これも不可能である。\path{computations/verify_cofactor_eighteen_quadric.py} は八点の積曲面での
長さ十二と両変形の交点数を再現する。
\end{proof}
